\begin{description}
\item[(T07.1)] The accretion radius can be estimated by equating the
  gravitational potential energy of the stellar wind to its kinetic energy
  \begin{align*}
    \frac{1}{2}\rho_{\mathrm{w}} v_{\mathrm{rel}}^2
      &= \frac{G M_{\mathrm{NS}} \rho_{\mathrm{w}}}{R_{\mathrm{acc}}}
      \tag*{(0.5)}\\
    \implies R_{\mathrm{acc}} &= \frac{2 G M_{\mathrm{NS}}}{v_{\mathrm{rel}}^2}
      \tag*{(0.5)}
  \end{align*}
  where $\rho_{\mathrm{w}}$ is the wind density and $v_{\mathrm{rel}}$ is the
  relative velocity of the wind.

  For the case considered above,
  \[
    v_{\mathrm{rel}}^2 = v_{\mathrm{w}}^2 + v_{\mathrm{orb}}^2
  \]
  Here, $v_{\mathrm{w}} \gg v_{\mathrm{orb}}$, therefore,
  $v_{\mathrm{rel}}^2 \approx v_{\mathrm{w}}^2
  = 9 \times 10^{12}\,\mathrm{m^2\,s^{-2}}$. \hfill(1.0)

  which gives,
  \[
    R_{\mathrm{acc}} = \frac{2 \times 6.674 \times 10^{-11} \times 2.0 \times
      1.988 \times 10^{30}}{9 \times 10^{12}}\,\mathrm{m}
  \]
  \[
    R_{\mathrm{acc}} = 5.9 \times 10^{4}\,\mathrm{km}
  \]
  \hfill(1.0)

\item[(T07.2)] In the case of spherically symmetric mass loss from the
  companion and $v_{\mathrm{w}} \gg v_{\mathrm{orb}}$, the companion mass loss
  is
  \[
    \dot{M}_{\mathrm{OB}} = 4\pi a^2 \rho_{\mathrm{w}} v_{\mathrm{w}}
  \]
  \hfill(1.0)

  where $a = 0.5\,\mathrm{au} = 7.480 \times 10^{10}$\,m is the orbital
  separation. Mass accretion rate on the compact object is
  \[
    \dot{M}_{\mathrm{acc}} \approx \pi R_{\mathrm{acc}}^2\,\rho_{\mathrm{w}}\,
      v_{\mathrm{rel}}
  \]
  \hfill(1.0)

  Also here, since $v_{\mathrm{w}} \gg v_{\mathrm{orb}}$,
  $v_{\mathrm{rel}} \approx v_{\mathrm{w}}$. This implies,
  \begin{align*}
    \dot{M}_{\mathrm{acc}}
      &\approx \left(\frac{\dot{M}_{\mathrm{OB}}}{4}\right)
        \left(\frac{R_{\mathrm{acc}}}{a}\right)^2
        \approx 9.74 \times 10^{11}\,\mathrm{kg\,s^{-1}}\\
    \implies \dot{M}_{\mathrm{acc}}
      &\approx 1.6 \times 10^{-11}\,\mathrm{M_\odot\,yr^{-1}}
      \tag*{(1.0)}
  \end{align*}

\item[(T07.3)] The orbital speed of the NS in the frame corresponding to the
  center of OB star (assuming the star is stationary) is given by
  \[
    v_{\mathrm{orb}}^2 = \frac{G M_{\mathrm{OB}}}{a}
  \]
  \hfill(1.0)

  Using vector algebra we get
  $\tan\beta = v_{\mathrm{orb}}/v_{\mathrm{w}}$. \hfill(1.0)

  Now, the ratio of the mass accretion rate to that of mass loss rate from the
  star
  \begin{align*}
    \frac{\dot{M}_{\mathrm{acc}}}{\dot{M}_{\mathrm{OB}}}
      &\sim \frac{\pi R_{\mathrm{acc}}^2 \rho_{\mathrm{w}} v_{\mathrm{rel}}}
        {4\pi a^2 \rho_{\mathrm{w}} v_{\mathrm{w}}}\\
      &= \frac{1}{4}\left(\frac{R_{\mathrm{acc}}}{a}\right)^2
        \sqrt{\frac{v_{\mathrm{orb}}^2 + v_{\mathrm{w}}^2}{v_{\mathrm{w}}^2}}
       = \frac{1}{4}\left(\frac{R_{\mathrm{acc}}}{a}\right)^2
        \sqrt{1 + \tan^2\beta}
      \tag*{(1.0)}
  \end{align*}
  To find the expression of the ratio of the two radii, $R_{\mathrm{acc}}$ and
  orbital separation $a$,
  \begin{align*}
    \frac{R_{\mathrm{acc}}}{a}
      &= \frac{2 G M_{\mathrm{NS}}}{v_{\mathrm{rel}}^2}\frac{1}{a}
       = \frac{2 G M_{\mathrm{OB}}\,q}{v_{\mathrm{rel}}^2}\frac{1}{a}
       = \frac{2 v_{\mathrm{orb}}^2}{v_{\mathrm{rel}}^2}\,q
      \tag*{(1.0)}\\
      &= \left(\frac{2\tan^2\beta}{1 + \tan^2\beta}\right) q
      \tag*{(1.0)}
  \end{align*}
  On substituting and simplifying, we can express,
  \[
    \implies f(\tan\beta, q)
      = q^2 \left(\frac{\tan^4\beta}{(1 + \tan^2\beta)^{3/2}}\right)
  \]
  \hfill(1.0)

\item[(T07.4a)] The magnetic pressure $P_{\mathrm{mag}} = B^2/2\mu_0$
  increases rapidly towards the NS magnetic poles.
  \[
    \implies P_{\mathrm{eq,\,mag}}
      = \frac{B_0^2 R_{\mathrm{NS}}^6}{2\mu_0 r^6}
  \]
  \hfill(1.0)

\item[(T07.4b)] The critical magnetic field $B_{0,\mathrm{c}}$ can be obtained
  by equating pressure due to incoming wind with the $P_{\mathrm{eq,\,mag}}$ at
  $r = R_{\mathrm{acc}}$. \hfill(1.0)

  The pressure due to incoming wind at $r = R_{\mathrm{acc}}$ is given by
  \[
    \rho_{\mathrm{w}} v_{\mathrm{rel}}^2 \sim \rho_{\mathrm{w}} v_{\mathrm{w}}^2
  \]
  \hfill(1.0)

  From part (T07.2), one can obtain the expression for $\rho_{\mathrm{w}}$ as
  follows:
  \[
    \rho_{\mathrm{w}}
      = \frac{\dot{M}_{\mathrm{acc}}}{\pi R_{\mathrm{acc}}^2 v_{\mathrm{w}}}
      = \frac{\dot{M}_{\mathrm{OB}}}{4\pi a^2 v_{\mathrm{w}}}
  \]
  \hfill(1.0)

  Now equating the pressures due to wind and magnetic field, we have:
  \[
    \frac{B_{0,\mathrm{c}}^2 R_{\mathrm{NS}}^6}{2\mu_0 R_{\mathrm{acc}}^6}
      = \frac{\dot{M}_{\mathrm{acc}}}{\pi R_{\mathrm{acc}}^2}
        \left(\frac{2 G M_{\mathrm{NS}}}{R_{\mathrm{acc}}}\right)^{1/2}
  \]
  \hfill(1.0)

  which gives the value of the critical magnetic field as
  \[
    \implies B_{0,\mathrm{c}}
      = \left(\frac{2\mu_0 \dot{M}_{\mathrm{acc}}
          \sqrt{2 G M_{\mathrm{NS}}}}{\pi R_{\mathrm{NS}}^6}\right)^{1/2}
        R_{\mathrm{acc}}^{7/4}
  \]
  \hfill(1.0)

  Now putting the values we get
  \[
    \implies B_{0,\mathrm{c}} \approx 4.0 \times 10^{9}\,\mathrm{T}
  \]
  \hfill(2.0)
\end{description}
